% Hafta 1 — Saltzer ve Schroeder (1975): güvenli tasarımın omurgası
% Derleme: py -3.12 tools/latex/derle.py docs/week-1/assets/h01-18-saltzer-schroeder.tex
\documentclass[tikz,border=8pt]{standalone}
\usepackage{cen429-cizim}
\begin{document}
\begin{tikzpicture}[node distance=10mm and 10mm]
  \node[baslik] (b) at (0,0) {Saltzer ve Schroeder (1975): güvenli tasarımın omurgası};
  \node[altbaslik] (alt) at (b.south west) {elli yıl sonra hâlâ geçerli sekiz ilke};

  \node[kutu=68mm, below=8mm of alt.south west, anchor=north west, xshift=2mm, text width=60mm] (i1)
    {\kb{En az ayrıcalık}\\işini yapacak kadar yetki};
  \node[kutu=68mm, right=of i1, text width=60mm] (i2)
    {\kb{Güvenli varsayılan}\\varsayılan: izin YOK};
  \node[kutu=68mm, below=of i1, text width=60mm] (i3)
    {\kb{Tam aracılık}\\her erişim her seferinde denetlenir};
  \node[kutu=68mm, right=of i3, text width=60mm] (i4)
    {\kb{Açık tasarım}\\gizlilik anahtarda, tasarımda değil};
  \node[kutu=68mm, below=of i3, text width=60mm] (i5)
    {\kb{Mekanizma ekonomisi}\\basit tasarım = az hata};
  \node[kutu=68mm, right=of i5, text width=60mm] (i6)
    {\kb{Ortak mekanizmayı azalt}\\paylaşılan yol = paylaşılan risk};
  \node[kutu=68mm, below=of i5, text width=60mm] (i7)
    {\kb{Yetki ayrımı}\\tek onay yetmez};
  \node[kutu=68mm, right=of i7, text width=60mm] (i8)
    {\kb{Kullanılabilirlik}\\zor gelen güvenlik atlanır};

  \node[dip, text width=155mm, below=9mm of i7.south -| i5, anchor=north]
    {Bu ilkeler bir denetim listesi değil, tasarım kararlarını tartarken kullanılan ölçülerdir.};
\end{tikzpicture}
\end{document}
